\documentclass[12pt,a4paper,reqno]{amsart}

% language
\usepackage[greek,english]{babel}
\usepackage[utf8]{inputenc}
\usepackage{alphabeta}

% change default names to greek
\addto\captionsenglish{
    \renewcommand{\contentsname}{Περιεχόμενα}
    \renewcommand{\refname}{Βιβλιογραφία}
    \renewcommand{\datename}{Ημερομηνία:}
    \renewcommand{\urladdrname}{Ιστοσελίδα}
}

% math 
\usepackage{amsmath,amsthm,amssymb,amscd,stmaryrd}

% font
\usepackage[cal=euler]{mathalfa}
\usepackage{libertinus-type1}
% \usepackage{txfonts} % for upright greek letters
\usepackage{bm} % for bold symbols
\usepackage{bbm} % for the simply-looking bb symbols

% miscellaneous 
\usepackage{changepage} %for indenting environments
\usepackage{csquotes} % example: \textcquote{}
\usepackage{blkarray}
\setcounter{MaxMatrixCols}{20} % default for pmatrix is 10!!
\usepackage{ytableau}
\usepackage{array} %needed to increase the vertical length in a tabular

% drawing
\usepackage{tikz,tikz-cd}
\usetikzlibrary{shapes.misc, patterns, matrix, calc, intersections,positioning,arrows.meta}
\usepackage{graphics,graphicx}
\usepackage{float} % provides enhanced control and customization options for floating objects such as figures and tables

% colors
\usepackage{xcolor}
\definecolor{darkcandyapplered}{rgb}{0.64, 0.0, 0.0}
\definecolor{midnightblue}{rgb}{0.1, 0.1, 0.44}
\definecolor{mylightblue}{HTML}{336699}
\definecolor{burntorange}{rgb}{0.8, 0.33, 0.0}
\definecolor{iceberg}{rgb}{0.44, 0.65, 0.82}
\definecolor{applegreen}{rgb}{0.55, 0.71, 0.0}
\definecolor{canaryyellow}{rgb}{1.0, 0.94, 0.0}
\definecolor{lightgray}{rgb}{0.83, 0.83, 0.83}

% hrefs
\usepackage{hyperref}
\usepackage[noabbrev,capitalize]{cleveref}
\hypersetup{
    pdftoolbar=true,        
    pdfmenubar=true,        
    pdffitwindow=false,     
    pdfstartview={FitH},  % fits the width of the page to the window
    pdftitle={},
    pdfauthor={},
    pdfsubject={},
    pdfkeywords={},
    pdfnewwindow=true,  % links in new window
    colorlinks=true,  % false: boxed links; true: colored links
    linkcolor=darkcandyapplered,   % color of internal links
    citecolor=midnightblue,  % color of links to bibliography
    urlcolor=cyan,  % color of external links
    linktocpage=true  % changes the links from the section body to the page number
    }

% geometry
\textwidth=16cm 
\textheight=21cm 
\hoffset=-55pt 
\footskip=25pt

% thm envs (you might need to change the path)
\theoremstyle{definition}
\newtheorem{theorem}{Θεώρημα}
\newtheorem{proposition}[theorem]{Πρόταση}
\newtheorem{lemma}[theorem]{Λήμμα}
\newtheorem{corollary}[theorem]{Πόρισμα}
\newtheorem{conjecture}[theorem]{Εικασία}
\newtheorem{problem}[theorem]{Πρόβλημα}
\newtheorem*{claim}{Ισχυρισμός}
\newtheorem{observation}[theorem]{Παρατήρηση}
\newtheorem{definition}[theorem]{Ορισμός}
\newtheorem{question}[theorem]{Ερώτημα}
\newtheorem*{questions}{Ερωτήματα}
\newtheorem*{que}{Ερώτημα}
\newtheorem{example}[theorem]{Παράδειγμα}
\newtheorem{exercise}{Άσκηση}
\newtheorem*{zorn}{Λήμμα του Zorn}
\newtheorem*{example_6.6}{Παράδειγμα~6.6}

\theoremstyle{remark}
\newtheorem*{remark}{Παρατήρηση}

% fixes the correct numbering of environments
\numberwithin{theorem}{section}
\numberwithin{exercise}{section}
\numberwithin{equation}{section}

% math ops 
% fields
\renewcommand{\AA}{\mathbbmss{A}}
\newcommand{\NN}{\mathbbmss{N}} 
\newcommand{\ZZ}{\mathbbmss{Z}} 
\newcommand{\QQ}{\mathbbmss{Q}} 
\newcommand{\RR}{\mathbbmss{R}} 
\newcommand{\CC}{\mathbbmss{C}} 
\newcommand{\KK}{\mathbbmss{K}} 
\newcommand{\FF}{\mathbbmss{F}} 
\newcommand{\HH}{\mathbbmss{H}} 
\newcommand{\VV}{\mathbbmss{V}} 
\newcommand{\II}{\mathbbmss{I}} 

% symmetric group
\newcommand{\fS}{\mathfrak{S}}  
\newcommand{\fm}{\mathfrak{m}}  

% calligraphic 
\newcommand{\aA}{\mathcal{A}} 
\newcommand{\bB}{\mathcal{B}}
\newcommand{\cC}{\mathcal{C}}
\newcommand{\dD}{\mathcal{D}}
\newcommand{\eE}{\mathcal{E}}
\newcommand{\fF}{\mathcal{F}}
\newcommand{\hH}{\mathcal{H}}
\newcommand{\iI}{\mathcal{I}}
\newcommand{\lL}{\mathcal{L}}
\newcommand{\oO}{\mathcal{O}}
\newcommand{\pP}{\mathcal{P}}
\newcommand{\qQ}{\mathcal{Q}}
\newcommand{\sS}{\mathcal{S}}
\newcommand{\mM}{\mathcal{M}}
\newcommand{\uU}{\mathcal{U}}
\newcommand{\zZ}{\mathcal{Z}}
\newcommand{\vV}{\mathcal{V}}

% bold
\newcommand{\bfa}{\pmb{a}}
\newcommand{\bfb}{\pmb{b}}
\newcommand{\bfe}{\mathbf{e}}
\newcommand{\bfF}{\pmb{F}}
\newcommand{\bfR}{\pmb{R}}
\newcommand{\bfv}{\mathbf{v}}
%\newcommand{\bfx}{\bm{x}}
%\newcommand{\bfx}{\mathbf{x}} 
\newcommand{\bfx}{\pmb{x}}
\newcommand{\bfX}{\pmb{X}}
\newcommand{\bfy}{\pmb{y}}
\newcommand{\bfz}{\pmb{z}}

% roman
\newcommand{\rmA}{\mathrm{A}}
\newcommand{\rmB}{\mathrm{B}}
\newcommand{\rmC}{\mathrm{C}}
\newcommand{\rmD}{\mathrm{D}} 
\newcommand{\rmF}{\mathrm{F}} 
\newcommand{\rmH}{\mathrm{H}} 
\newcommand{\rmh}{\mathrm{h}} 
\newcommand{\rmI}{\mathrm{I}} 
\newcommand{\rmK}{\mathrm{K}}
\newcommand{\rmM}{\mathrm{M}}
\newcommand{\rmP}{\mathrm{P}}  
\newcommand{\rmp}{\mathrm{p}}  
\newcommand{\rmQ}{\mathrm{Q}}  
\newcommand{\rmR}{\mathrm{R}}
\newcommand{\rmS}{\mathrm{S}}
\newcommand{\rmT}{\mathrm{T}}
\newcommand{\rmU}{\mathrm{U}}
\newcommand{\rmV}{\mathrm{V}}
\newcommand{\rmY}{\mathrm{Y}}
\newcommand{\rmZ}{\mathrm{Z}}
\newcommand{\rmz}{\mathrm{z}}

% arrows and symbols 
\renewcommand{\to}{\rightarrow}
\newcommand{\toto}{\longrightarrow}
\newcommand{\mapstoto}{\longmapsto}
\newcommand{\then}{\Rightarrow}
\newcommand{\IFF}{\Leftrightarrow}
\newcommand{\tl}{\tilde}
\newcommand{\wtl}{\widetilde}
\newcommand{\ol}{\overline}
\newcommand{\ul}{\underline}
\newcommand{\oldemptyset}{\emptyset}
\renewcommand{\emptyset}{\varnothing}
\DeclareMathSymbol{\Arg}{\mathbin}{AMSa}{"39} % for arguments 
\newcommand{\onto}{\ensuremath{\twoheadrightarrow}}
\newcommand{\tle}{\trianglelefteq}
\newcommand{\tge}{\trianglerighteq}

% custom ops
\newcommand{\rmKer}{\mathrm{Ker}}
\newcommand{\rmIm}{\mathrm{Im}}
\newcommand{\lcm}{\mathrm{lcm}}
\newcommand{\GL}{\mathrm{GL}}
\newcommand{\ev}{\mathrm{ev}}
\newcommand{\End}{\mathrm{End}}
\newcommand{\rmid}{\mathrm{id}}
\newcommand{\rmspan}{\mathrm{span}}
\newcommand{\Hom}{\mathrm{Hom}}
\newcommand{\Ann}{\mathrm{Ann}}
\newcommand{\Set}{\pmb{\mathrm{Set}}}
\newcommand{\Rmod}{\pmb{\mathrm{mod}}}
\newcommand{\rmrank}{\mathrm{rank}}
\newcommand{\mSpec}{\mathrm{mSpec}}
\newcommand{\Spec}{\mathrm{Spec}}
\newcommand{\tor}{\mathrm{tor}}
\newcommand{\Pol}[1]{F[x_1, x_2, \dots, x_{#1}]}
\newcommand{\rad}{\mathrm{rad}}

% absolute value symbol
\usepackage{mathtools} 
\DeclarePairedDelimiter\abs{\lvert}{\rvert}%
\DeclarePairedDelimiter\norm{\lVert}{\rVert}%
\makeatletter
\let\oldabs\abs
\def\abs{\@ifstar{\oldabs}{\oldabs*}}

% tensor symbol
\newcommand{\tensor}[1]{%
  \mathbin{\mathop{\otimes}\limits_{#1}}%
}

% permutation cycle notation
\ExplSyntaxOn
\NewDocumentCommand{\cycle}{ O{\;} m }
 {
  (
  \alec_cycle:nn { #1 } { #2 }
  )
 }

\seq_new:N \l_alec_cycle_seq
\cs_new_protected:Npn \alec_cycle:nn #1 #2
 {
  \seq_set_split:Nnn \l_alec_cycle_seq { , } { #2 }
  \seq_use:Nn \l_alec_cycle_seq { #1 }
 }
\ExplSyntaxOff

% setminus symbol
\newcommand{\mysetminusD}{\hbox{\tikz{\draw[line width=0.6pt,line cap=round] (3pt,0) -- (0,6pt);}}}
\newcommand{\mysetminusT}{\mysetminusD}
\newcommand{\mysetminusS}{\hbox{\tikz{\draw[line width=0.45pt,line cap=round] (2pt,0) -- (0,4pt);}}}
\newcommand{\mysetminusSS}{\hbox{\tikz{\draw[line width=0.4pt,line cap=round] (1.5pt,0) -- (0,3pt);}}}
\newcommand{\sm}{\mathbin{\mathchoice{\mysetminusD}{\mysetminusT}{\mysetminusS}{\mysetminusSS}}}

% miscellaneous commands
\newcommand{\defn}[1]{{\color{mylightblue}{#1}}}
\newcommand{\toDo}{{\bf\color{red} TODO}}
\newcommand{\toCite}{{\bf\color{green} CITE}}
\newcommand*{\vertbar}{\rule[-1ex]{0.5pt}{2.5ex}} % for matrices with column vectors
\newcommand*{\horzbar}{\rule[.5ex]{2.5ex}{0.5pt}} % for matrices with row vectors
\newcommand{\myblue}[1]{{\color{iceberg}{#1}}}
\newcommand{\myorange}[1]{{\color{burntorange}{#1}}}
\newcommand{\mygreen}[1]{{\color{applegreen}{#1}}}
\newcommand{\myred}[1]{{\color{darkcandyapplered}{#1}}}

% ferrer's diagram
\newcommand{\fdiagram}[1]{
    \begin{tikzpicture}[scale=.7]
        \fill foreach \Z [count=\Y] in {#1}
        {foreach \X in {1,...,\Z} 
        {(\X,-\Y) circle[radius=3pt]}};
    \end{tikzpicture}
}

%
\usepackage{pgfplots}
\pgfplotsset{compat=1.18}

%
\makeatletter
\def\mathcolor#1#{\@mathcolor{#1}}
\def\@mathcolor#1#2#3{%
  \protect\leavevmode
  \begingroup
    \color#1{#2}#3%
  \endgroup
}
\makeatother

% restriction
\newcommand\restr[2]{{% we make the whole thing an ordinary symbol
  \left.\kern-\nulldelimiterspace % automatically resize the bar with \right
  #1 % the function
  \littletaller % pretend it's a little taller at normal size
  \right|_{#2} % this is the delimiter
  }}
\newcommand{\littletaller}{\mathchoice{\vphantom{\big|}}{}{}{}}

% 
\newenvironment{nouppercase}{%
  \let\uppercase\relax%
  \renewcommand{\uppercasenonmath}[1]{}}{}

% titlepage
\title{226: Θεωρία Δακτυλίων και Modules}
\author[Β.~Δ. Μουστακας]{Βασίλης Διονύσης Μουστάκας \\ Πανεπιστήμιο Κρήτης}
\date{7 Μαΐου 2026}
% \urladdr{\href{https://sites.google.com/view/vasmous}{https://sites.google.com/view/vasmous}}

\begin{document}

\begingroup
\def\uppercasenonmath#1{} % this disables uppercase title
\let\MakeUppercase\relax % this disables uppercase authors
\maketitle
\endgroup

\setcounter{section}{9}
\setcounter{theorem}{6}
\setcounter{equation}{2}
\begin{center}
    \textbf{9. Αλγεβρική Γεωμετρία: Εισαγωγή
} (Συνέχεια)
\end{center}


\begin{example}
    \label{ex:vanishing_ideal_of_whole_variety}
    Υπενθυμίζουμε ότι το $F$ είναι ένα άπειρο σώμα. Θα δείξουμε ότι 
    \[
    \iI(\AA_F^n) = 0.
    \]
    Με άλλα λόγια\footnote{Ταυτίζοντας τα πολυώνυμα στις $n$ μεταβλητές με τις πολυωνυμικές συναρτήσεις $\AA^n \to \AA^1$, όπως κάναμε στο τέλος της Ενότητας 3, η παρατήρηση αυτή μας πληροφορεί κάτι γεωμετρικό χρησιμοποιώντας άλγεβρα.}, κάθε μη μηδενικό πολυώνυμο σε $n$ μεταβλητές είναι μη μηδενικό σε κάποιο σημείο του $\AA^n$ ή ισοδύναμα αν ένα πολυώνυμο μηδενίζεται σε ολόκληρο το $\AA^n$, τότε πρέπει αναγκαστικά να είναι το μηδενικό πολυώνυμο. Θα κάνουμε επαγωγή ως προς $n$. 

    Για $n=1$, επειδή το πλήθος των ριζών ενός μη μηδενικού πολυωνύμου $f \in F[x_1]$ είναι το πολύ $\deg(f)$ (βλ. συζήτηση μετά την Πρόταση 3.7), έπεται ότι $f \notin \iI(\AA^1)$, διότι το $F$ έχει άπειρα στοιχεία. Άρα, $\iI(\AA^1) = 0$. Αυτή είναι η βάση της επαγωγής. 
    
    Υποθέτουμε ότι $n > 1$, $\iI(\AA^{n-1}) = 0$ και θεωρούμε ένα μη μηδενικό $f \in \Pol{n}$. Αρκεί να δείξουμε ότι $f \notin \iI(\AA^n)$. Γράφουμε 
    \begin{equation}
        \label{eq:f}
        f = p_0(x_1,x_2,\dots, x_{n-1}) + p_1(x_1,x_2,\dots, x_{n-1}) x_n + \cdots + p_d(x_1,x_2,\dots, x_{n-1}) x_n^d,
    \end{equation}
    για κάποια $p_1, p_2, \dots, p_d \in \Pol{n-1}$ και κάποιο $d \in \NN$ με $p_d \neq 0$. Από την επαγωγική υπόθεση, υπάρχει ένα $(a_1, a_2, \dots, a_{n-1}) \in \AA^{n-1}$, τέτοιο ώστε $p_d(a_1, a_2, \dots, a_{n-1}) \neq 0$. Αντικαθιστώντας στην \eqref{eq:f}, έπεται ότι το πολυώνυμο
    \[
    f(a_1, a_2, \dots, a_{n-1}, x_n) \ \in F[x_n]
    \] 
    είναι μη μηδενικό. Από την βάση της επαγωγής (η περίπτωση της μιας μεταβλτητής), υπάρχει $a_n \in \AA^1$ τέτοιο ώστε 
    \[     
    f(a_1, a_2, \dots, a_n) \neq 0. 
    \]   
    Επομένως, $f \notin \iI(\AA^n)$ και η επαγωγή ολοκληρώνεται.
\end{example}

\begin{example}
    \label{ex:vanishing_ideal_of_a_point}
    Έστω $\bfa = (a_1, a_2, \dots, a_n) \in \AA^n$. Στο τέλος της Ενότητας 3 αναφέραμε ότι 
    \[
    \fm_{\bfa} \coloneqq \left(
        x_1-a_1, x_2-a_2, \dots, x_n-a_n
    \right) = \ker\left(\ev_{\bfa}\right).
    \]
    Συνεπώς, $f \in \fm_{\bfa}$ αν και μόνο αν $f(\bfa) = 0$ ή ισοδύναμα, αν και μόνο αν $f \in \iI(\{\bfa\})$. Άρα, 
    \[
    \iI(\{\bfa\}) = \fm_{\bfa},
    \]
    τ' οποίο είναι μεγιστικό ιδεώδες του $\Pol{n}$ (γιατί;).

    % Πράγματι, αν $f \in \iI(\{\bfa\})$, τότε $f(\bfa) = f(a_1, a_2, \dots, a_n) = 0$. Αυτό σημαίνει ότι $f \in (x_i-a_i)$, για κάθε $1 \le i \le n$ και γι αυτό $f \in \left(
    %     x_1-a_1, x_2-a_2, \dots, x_n-a_n
    % \right)$. Επομένως, $\iI(\{\bfa\}) \subseteq \fm_\bfa$. Από την άλλη μεριά, αν $f \in \left(
    %     x_1-a_1, x_2-a_2, \dots, x_n-a_n
    % \right)$, τότε 
    % \[
    % f(x_1, x_2, \dots, x_n) = \sum_{i=1}^n g_i(x_1,x_2, \dots, x_n)(x_i-a_i)
    % \]
    % για κάποια $g_1, g_2, \dots, g_n \in \Pol{n}$. Συνεπώς, $f(\bfa) = 0$ και γι αυτό $\fm_\bfa \subseteq \iI(\{\bfa\})$.

    Η ισότητα $\fm_{\bfa} = \ker(\ev_{\bfa})$ δεν είναι απολύτως προφανής. Μια αυστηρή απόδειξη μπορεί να γίνει με τρόπο όμοιο με αυτό του Παραδείγματος \ref{ex:vanishing_ideal_of_whole_variety}. Για την κατεύθυνση \textquote{$\subseteq$}, αν 
    \[
    f = f_1(x_1-a_1) + f_2(x_2-a_2) + \cdots + f_n(x_n-a_n),
    \]
    για κάποια $f_1, f_2, \dots, f_n \in \Pol{n}$, τότε $\ev_{\bfa}(f) = 0$. Θα δείξουμε την άλλη κατεύθυνση, δηλαδή ότι $\ker(\ev_{\bfa}) \subseteq \fm_{\bfa}$ με επαγωγή στο $n$. Για $n=1$, το ζητούμενο έπεται από την Πρόταση 3.7 (γιατί;). Αυτή είναι η βάση της επαγωγής. 
    
    Έστω $n > 1$, υποθέτουμε ότι 
    \[
    \ker(\ev_{\bfb}) \subseteq \fm_{\bfb}
    \] 
    για κάθε $\bfb \in \AA^{n-1}$ και θεωρούμε ένα $f \in \ker(\ev_{\bfa})$. Γράφουμε το $f$ όπως στην Εξίσωση \eqref{eq:f}. Διακρίνουμε περιπτώσεις. Αν $d=0$, τότε 
    \[
    f = p_0(x_1,x_2,\dots, x_{n-1}) \ \in \ \Pol{n-1}
    \]
    και γι αυτό $\ev_{\hat{\bfa}}(f) = 0$, όπου $\hat{\bfa} = (a_1, a_2, \dots, a_{n-1})$. Από την επαγωγική υπόθεση έπεται ότι 
    \[
    f \in \fm_{\hat{\bfa}} \subseteq \fm_{\bfa}.
    \]

    Τώρα, υποθέτουμε ότι $d > 1$. Εφαρμόζοντας την γενικευμένη εκδοχή του Θεωρήματος 3.3 (βλ. συζήτηση αμέσως μετά την διατύπωση του θεωρήματος) για την ακέραια περιοχή $R = \Pol{n-1}$ και $g(x_n) = x_n-a_n$, έπεται ότι 
    \[
    f = q(x_n-a_n) + r,
    \]
    για κάποια $q, r \in R[x_n]$ με το $r$ να έχει βαθμό μικρότερο από $1$ ως πολυώνυμο στο $x_n$. Αυτό σημαίνει ότι
    \[  
    r \in R = \Pol{n-1}
    \] 
    και επειδή $f \in \ker(\ev_{\bfa})$, έπεται ότι $\ev_{\hat{\bfa}}(r) = 0$ (γιατί;), δηλαδή $r \in \ker(\ev_{\hat{\bfa}})$. Από την επαγωγική υπόθεση, αυτό έχει ως συνέπεια $r \in \fm_{\hat{\bfa}}$, δηλαδή 
    \[
    r = f_1(x_1-a_1) + f_2(x_2-a_2) + \cdots + f_{n-1}(x_{n-1}-a_{n-1})
    \] 
    για κάποια $f_1, f_2, \dots, f_{n-1} \in \Pol{n-1}$ και γι αυτό 
    \[
    f = f_1(x_1-a_1) + f_2(x_2-a_2) + \cdots + f_{n-1}(x_{n-1}-a_{n-1}) + q(x_n-a_n).
    \]
    Άρα, $f \in \fm_{\bfa}$ και το ζητούμενο έπεται.
\end{example}

\begin{proposition}
    \label{prop:vanishing_ideal}
    Έστω $I$ ένα ιδεώδες του $\Pol{n}$ και $A \subseteq \AA^n$. Τότε 
    \begin{equation} 
        \label{eq:one_way}
        A \subseteq \vV\left(
        \iI(A)
        \right)
    \end{equation}
    με ισότητα να ισχύει αν και μόνο αν το $A$ είναι πολλαπλότητα, και 
    \begin{equation} 
        \label{eq:other_way}
        I \subseteq \iI(\vV(I)).
    \end{equation}
\end{proposition}

\begin{proof}[Απόδειξη]
    Για την \eqref{eq:one_way}, αν $\bfa \in A$, τότε $f(\bfa) = 0$, για κάθε $f \in \iI(A)$, εξ' ορισμού και κατά συνέπεια $\bfa \in \vV\left(
        \iI(A)
    \right)$. Επομένως, $A \subseteq \vV\left(
        \iI(A)
    \right)$. Ο δεύτερος εγκλεισμός αποδεικνύεται με παρόμοιο τρόπο.

    Για την ισότητα στην \eqref{eq:one_way}, αν το $A$ είναι της μορφής $\vV(S)$ για κάποιο $S \subseteq \Pol{n}$, τότε $S \subseteq \iI(\vV(S)) = \iI(A)$. Συνεπώς, 
    \[
    \vV\left(
        \iI(A)
    \right) \subseteq \vV(S) = A,
    \]
    από την παρατήρηση μετά τον Ορισμό 9.1 και το ζητούμενο έπεται.
\end{proof}

\begin{example}
\label{ex:whatever}
Οι ανισότητες της Πρότασης \ref{prop:vanishing_ideal} κάποιες φορές είναι αυστηρές. Για παράδειγμα, αν $A = \{(t, t^2) : t \in \NN\} \subseteq \AA_\RR^2$
\[
\begin{tikzpicture}
    \begin{axis}[
        axis lines = middle,
        axis line style={-latex, gray!50},
        xlabel = {$x$},
        ylabel = {$y$},
        label style = {gray!50},
        xmin=-2, xmax=2,
        ymin=-2, ymax=4,
        xticklabels=\empty,
        yticklabels=\empty,
        ticks=none
        ]
        \addplot[black!30] {x^2}; 

        \addplot[
            only marks, 
            mark=*, 
            mark size=1.5pt, 
            domain=0:2, 
            samples=20
        ] {x^2};

    \end{axis}
\end{tikzpicture}
\]
τότε $\iI(A) = (y-x^2)$ και γι αυτό 
\[
A \subsetneqq  \{(t,t^2) : t \in \RR\} = \vV(y-x^2) = \vV\left(
        \iI(A)
    \right).
\] 
Πράγματι, αν $f \in (y-x^2)$, τότε $f(x,y) = g(x,y)(y-x^2)$ για κάποιο $g \in \RR[x,y]$. Συνεπώς, 
\[
f(t, t^2) = g(t,t^2) (t^2-t^2) = 0
\]
για κάθε $t \in \NN$. Επομένως, $f \in \iI(A)$ και γι αυτό $(y-x^2) \subseteq \iI(A)$. Αντιστρόφως, αν $f \in \iI(A)$, όμοια με το Παραδείγματα \ref{ex:vanishing_ideal_of_whole_variety} και \ref{ex:vanishing_ideal_of_a_point}, βλέπουμε το $f$ ως στοιχείο του $\mathcolor{gray}{\RR[x]}[y]$ και γι αυτό 
\[
f = q(y-\mathcolor{gray}{x^2}) + r
\]
για κάποια $q \in \mathcolor{gray}{\RR[x]}[y]$ και $r \in \mathcolor{gray}{\RR[x]}$. Επειδή $f \in \iI(A)$, έπεται ότι 
\[
r(t) = f(t,t^2) - q(\underbrace{t^2-t^2}_{= \ 0}) = 0
\]
για κάθε $t \in \NN$. Συνεπώς, όμοια με το Παράδειγμα 9.2 (θ), έχουμε $r = 0$ και γι αυτό $f \in (y-x^2)$. Επομένως $\iI(A) \subseteq (y-x^2)$.

Για τη δεύτερη ανισότητα, αν $I = \left( (y-x^2)^2 \right)$, τότε 
\[
\vV(I) = \{(t, t^2) : t \in \RR\}
\]
(γιατί;). Όμοια με προηγουμένως, παρατηρεί κανείς ότι 
$
\iI\left(
    \vV(I)
\right) = (y-x^2)$ και γι αυτό
\[
I = \left( (y-x^2)^2 \right) \subsetneqq (y-x^2) = \iI(\vV(I)).
\]
\end{example}

Η ισότητα \eqref{eq:one_way} στην περίπτωση όπου το $A$ είναι πολλαπλότητα μας πληροφορεί ότι οι απεικονίσεις 
\begin{align*}
\left\{
\begin{aligned}
&\qquad \quad \text{ιδεώδη} \\
& I \subseteq \Pol{n}
\end{aligned}
\right\}
&\to
\left\{
\begin{aligned}
&\text{πολλαπλότητες} \\
&\quad \ \ V \subseteq \AA_F^n
\end{aligned}
\right\} \\
I &\mapsto \vV(I) \\
\iI(A) &\mapsfrom A,
\end{align*}
είναι επί και ένα-προς-ένα, αντίστοιχα (γιατί;). Πότε ισχύει η ισότητα στην \eqref{eq:other_way}; Με άλλα λόγια, πότε είναι οι απεικονίσεις $\vV(\cdot)$ και $\iI(\cdot)$ ένα-προς-ένα και επί, αντίστοιχα;

Ας θυμηθούμε την έννοια του ριζικού ιδεώδους από την Άσκηση 12.

\begin{definition}
    \label{def:radical_ideal}
    Αν $I$ είναι ιδεώδες του $\Pol{n}$, τότε το σύνολο\footnote{Σε κάποια συγγράματα συμβολίζεται και ως $\rad(I)$.} 
    \[
    \sqrt{I} \coloneqq \{f \in \Pol{n} : f^k \in I, \ \text{για κάποιο $k \ge 1$}\}
    \]
     ονομάζεται \defn{ριζικό} (radical) του $I$. Το $I$ ονομάζεται \defn{ριζικό ιδεώδες} (radical ideal) αν $\sqrt{I} = I$.
\end{definition}

Στην Άσκηση 11 είδαμε ότι το ριζικό ενός ιδεώδους είναι ιδεώδες το οποίο το περιέχει.

\begin{proposition}
    \label{prop:vanishing_ideal_is_radical}
    Το ιδεώδες μηδενισμού ενός υποσυνόλου του αφφινικού χώρου είναι ριζικό ιδεώδες.
\end{proposition}

\begin{proof}[Απόδειξη]
    Έστω $A \subseteq \AA^n$. Αρκεί να δείξουμε ότι $\sqrt{\iI(A)} \subseteq \iI(A)$. Πράγματι, αν $f \in \sqrt{\iI(A)}$, τότε 
    \[
    0 = f^k(\bfa) = \left(f(\bfa)\right)^k,
    \] 
    για κάθε $\bfa \in A$. Συνεπώς, $f(\bfa) = 0$, για κάθε $\bfa \in A$ και γι αυτό $f \in \iI(A)$.
\end{proof}

Η Πρόταση \ref{prop:vanishing_ideal_is_radical} μας πληροφορεί ότι αυτό που \textquote{εμποδίζει} την ισότητα στην \eqref{eq:other_way} είναι το κατά πόσο το $I$ είναι ριζικό ιδεώδες. Στο παράδειγμα, το $I = \left( (y-x^2)^2 \right)$ δεν είναι ριζικό, καθώς $\sqrt{I} = (y-x^2)$ (γιατί;). Με άλλα λόγια, στο παράδειγμα αυτό ισχύει ότι
\[
\iI(\vV(I)) = \sqrt{I}.
\]
Αν ίσχυε πάντα αυτή η ταυτότητα, τότε θα είχαμε ισότητα στην \eqref{eq:other_way} αν και μόνο αν το $I$ ήταν ριζικό ιδεώδες. Ισχύει όμως αυτό σε κάθε περίπτωση; 


% \begin{thebibliography}{99}
%     % \bibitem{DF04}
%     % D. S. Dummit και R. M. Foote, 
%     % \textit{Abstract algebra},
%     % third edition, Wiley, 2004. 
%     % \bibitem{Mun14}
%     % J.~Munkres, 
%     % \textit{Topology},
%     % second edition, Prentice Hall, 2000.
% \end{thebibliography}
\end{document}